\clearpage
\fancyhead[L]{\small HUMBOLDT. CONSPECTUS LONGITUDINUM ETC.}
\hypertarget{native-514-lower}{}\label{native-514-lower}
\begin{center}\small Göttingische gelehrte Anzeigen. No. 193. p. 1928. 3 December 1808.\\[3pt]\rule{.72\linewidth}{.3pt}\end{center}

\emph{Conspectus longitudinum et latitudinum geographicarum, per decursum annorum} 1799 \emph{ad} 1804 \emph{in plaga aequinoctiali ab} \textsc{Alex. de Humboldt} \emph{astronomice observatarum, calculo subjecit} \textsc{Jabbo Oltmans}. 16 pages in folio. Available from \textsc{F. Schöll} and \textsc{J. G. Cotta} in Paris and Tübingen.

As is well known, a separate part of Mr \textsc{von Humboldt}’s great work is to be devoted to the valuable gains for the geography of tropical America from his travels. There we can expect to find all of his numerous astronomical observations in full, together with a very careful analysis of them by Mr \textsc{Jabbo Oltmans}, who has already earned a good reputation. Printing that work will take considerable time, however, so this summary compilation of the principal geographical positions has been prepared in the meantime. It contains only the final results. Nearly three hundred determinations are given here, based for the most part on Mr \textsc{von Humboldt}’s own observations. One cannot help admiring the energy and versatility of a man who took on so many different aims at once and knew how to achieve each as if he had been able to devote his time and energies to it alone.
